\documentclass[11pt,a4paper]{article}
\usepackage{gtu-el-yazisi}
\begin{document}
\notbilgi{Sayısal İşaret İşleme}{Frekans yanıtı, DTFT ve Z dönüşümü}{Dsp-ders-4.pdf}

\sayfa{1}
\elcizimi[0.55]{sekiller/Dsp-ders-4-s1-1.png}[1]

\[ \begin{aligned} H(e^{j\omega}) &= |H(e^{j\omega})|\,e^{j\theta} \\ \theta &= \angle H(e^{j\omega}) \end{aligned} \]

\[ \cos(\omega_0 n) = \tfrac12 e^{j\omega_0 n} + \tfrac12 e^{-j\omega_0 n} \]
\[ y[n] = \tfrac12 e^{j\omega_0 n} H(e^{j\omega_0}) + \tfrac12 e^{-j\omega_0 n} H(e^{-j\omega_0}) \]

\elcizimi[0.4]{sekiller/Dsp-ders-4-s1-2.png}[1]

\formulkutu[önemli]{\begin{aligned} h[n]\ \text{reel} \;\Rightarrow\; H(e^{j\omega}) &= H^{*}(e^{-j\omega}) \\ H(e^{-j\omega}) &= H^{*}(e^{j\omega}) \end{aligned}}

\[ \begin{aligned}
h[n]\ \text{reel} \;\Rightarrow\; y[n] &= \tfrac12 e^{j\omega_0 n} H(e^{j\omega_0}) + \tfrac12 e^{-j\omega_0 n} H^{*}(e^{j\omega_0}) \\
&= \tfrac12 e^{j\omega_0 n} |H(e^{j\omega_0})|\,e^{j\theta_0} + \tfrac12 e^{-j\omega_0 n} |H(e^{j\omega_0})|\,e^{-j\theta_0} \\
&= |H(e^{j\omega_0})| \left[ \frac{e^{j(\omega_0 n+\theta_0)} + e^{-j(\omega_0 n+\theta_0)}}{2} \right] \\
&= |H(e^{j\omega_0})| \cos(\omega_0 n+\theta_0)
\end{aligned} \]

\[ x[n] = \cos(0{,}1\pi n) + \cos(0{,}4\pi n) + \cos(0{,}6\pi n) \]
\[ y[n] = 0{,}8\cos(0{,}1\pi n+\theta_1) + 0{,}2\cos(0{,}4\pi n+\theta_2) + 0\cdot\cos(0{,}6\pi n+\theta_3) \]

\elcizimi[0.45]{sekiller/Dsp-ders-4-s1-3.png}[1]

\[ |H(e^{j\frac{\pi}{2}})|\cos\!\Big(\tfrac{\pi}{2}n + \angle H(e^{j\frac{\pi}{2}})\Big) \]

\[ x[n] = \cos\!\Big(\tfrac{\pi}{2}n\Big) \]
\[ h[n] = \delta[n] - \delta[n-1] \]
\begin{dikkat}
Sakın ha $y[n]=\cos\!\big(\tfrac{\pi}{2}n\big)-\cos\!\big(\tfrac{\pi}{2}(n-1)\big)$
\end{dikkat}
\[ y[n] = \,? \]

\[ \begin{aligned}
H(e^{j\omega}) &= 1 - e^{-j\omega} = 2j\,\frac{e^{-j\omega/2}\left(e^{j\omega/2}-e^{-j\omega/2}\right)}{2j} \\
&= j\,e^{-j\omega/2}\,2\sin(\omega/2)
\end{aligned} \]
\[ |H(e^{j\omega})| = 2\left|\sin(\omega/2)\right| \]
\[ \angle H(e^{j\omega}) = \begin{cases} \dfrac{\pi}{2} - \dfrac{\omega}{2}, & 0<\omega<\pi \\[4pt] \dfrac{\pi}{2} - \dfrac{\omega}{2} + \pi, & -\pi<\omega<0 \end{cases} \]

\elcizimi[0.45]{sekiller/Dsp-ders-4-s1-4.png}[1]

\sayfa{2}
\[ y[n] = \sqrt{2}\cos\!\Big(\tfrac{\pi}{2}n + \tfrac{\pi}{4}\Big) \]

\elcizimi[0.45]{sekiller/Dsp-ders-4-s2-1.png}[2]
\[ y[n] = \tfrac13 x[n] + \tfrac13 x[n-1] + \tfrac13 x[n-2] \]

\[ h[n] = \tfrac13\delta[n] + \tfrac13\delta[n-1] + \tfrac13\delta[n-2] \]

\[ \begin{aligned}
H(e^{j\omega}) &= \tfrac13 + \tfrac13 e^{-j\omega} + \tfrac13 e^{-j2\omega} \\
&= \sum_{k=0}^{2} \tfrac13 e^{-j\omega k}
\end{aligned} \]

\formulkutu[önemli]{\sum_{k=N_1}^{N_2} a^k = \frac{a^{N_1}-a^{N_2+1}}{1-a}}

\[ \begin{aligned}
&= \tfrac13\,\frac{(e^{-j\omega})^{0}-(e^{-j\omega})^{2+1}}{1-e^{-j\omega}} \\
&= \tfrac13\,\frac{1-e^{-j3\omega}}{1-e^{-j\omega}} \;\overset{\text{\scriptsize ikisini eşit bölse}}{=}\; \tfrac13\,\frac{e^{-j\frac32\omega}\left(e^{j\frac32\omega}-e^{-j\frac32\omega}\right)}{e^{-j\frac{\omega}{2}}\left(e^{j\frac{\omega}{2}}-e^{-j\frac{\omega}{2}}\right)} \\
&= \tfrac13\,e^{-j\omega}\,\frac{\sin\frac32\omega}{\sin\frac{\omega}{2}}
\end{aligned} \]

\begin{dikkat}
Her zaman sin hale getir
\end{dikkat}

\[ |H(e^{j\omega})| = \tfrac13\left|\frac{\sin\frac32\omega}{\sin\frac{\omega}{2}}\right| \]
\[ \angle H(e^{j\omega}) = -\omega + \pi \quad \text{(negatif olduğunda)} \]

\elcizimi[0.8]{sekiller/Dsp-ders-4-s2-2.png}[2]

\elcizimi[0.8]{sekiller/Dsp-ders-4-s2-3.png}[2]

\sayfa{3}
\elcizimi[0.6]{sekiller/Dsp-ders-4-s3-1.png}[3]

\elcizimi[0.4]{sekiller/Dsp-ders-4-s3-2.png}[3]
\[ y[n] = \begin{cases} 0, & n<0 \\ ?, & n\geq 0 \end{cases} \]

\elcizimi[0.6]{sekiller/Dsp-ders-4-s3-3.png}[3]

\[ y[n] = \sum_{k} h[k]\,x[n-k] \]
system is causal $\Rightarrow h[n]=0,\quad n<0$

\[ \begin{aligned}
&= \sum_{k=-\infty}^{\infty} h[k]\,e^{j\omega(n-k)}\,u[n-k] \\
&= \sum_{k=0}^{\infty} h[k]\,e^{j\omega(n-k)}\,u[n-k]
\end{aligned} \]
$h[n]$ casual

\formulkutu[önemli]{u[n-k]=0,\quad \begin{matrix} n-k<0 \\ k>n \end{matrix}}

\[ \begin{aligned}
y[n] &= \sum_{k=0}^{n} h[k]\,e^{j\omega(n-k)} \\
&= \sum_{k=0}^{\infty} h[k]\,e^{-j\omega k}\,e^{j\omega n} - \sum_{k=n+1}^{\infty} h[k]\,e^{-j\omega k}\,e^{j\omega n} \\
&= \underbrace{e^{j\omega n}H(e^{j\omega})}_{\text{steady state cevabı}} - \sum_{k=n+1}^{\infty} h[k]\,e^{-j\omega k}\,e^{j\omega n}
\end{aligned} \]

\[ y[n] = y_{ss}[n] + y_{tr}[n] \]

\[ \begin{aligned}
y[n] &= y_{ss}[n] + y_{tr}[n] \\
y_{ss}[n] &= e^{j\omega n}H(e^{j\omega}) \\
y_{tr}[n] &= -\sum_{k=n+1}^{\infty} h[k]\,e^{-j\omega k}\,e^{j\omega n}
\end{aligned} \]

\sayfa{4}
\elcizimi[0.7]{sekiller/Dsp-ders-4-s4-1.png}[4]

Filtrenin genişliği $m+1$

\[ \begin{aligned}
y_{tr}[0] &= -\sum_{k=1}^{\infty} (\ \ \ \ \ ) \neq 0 \\
y_{tr}[1] &= -\sum_{k=2}^{\infty} (\text{-- --}) \neq 0 \\
y_{tr}[2] &= -\sum_{k=3}^{\infty} (\text{-- -- --}) \neq 0 \\
&\ \ \vdots \\
y_{tr}[m-1] &= -\sum_{k=m}^{\infty} (\text{-- -- --}) \neq 0 \\
y_{tr}[m] &= \sum_{k=m+1}^{\infty} (\text{-- -- --}) = 0 \\
y_{tr}[m+1] &= \sum_{k=m+2}^{\infty} (\text{-- -- --}) = 0
\end{aligned} \]

pencere genişliği $-1$'e kadar 0'dan farklı;

pencere genişliği ne kadar dar olursa o kadar pencere genişliği dar olur

\[ h[n] = \tfrac13\delta[n] + \tfrac13\delta[n-1] + \tfrac13\delta[n-2] \]

$x[n]=1$ olsaydı, $y[n]=1$ olurdu

\[ x[n] = 1\cdot u[n] \]
\[ y[n] = 0, \quad n<0 \]
\[ \left.\begin{aligned} y[0] &= \tfrac13 \\ y[1] &= \tfrac23 \end{aligned}\right\} \text{Transient iki örnek \okunamadi{sürdü}} \]
\[ y[2] = 1 \]
\[ y[3] = 1 \]

\begin{dikkat}
pencere genişliği eksi 1'i kadar
\end{dikkat}

\vurgu[pembe]{Peki transient kararlı mı ??}

\[ \begin{aligned}
|y_{tr}[n]| &= \left|\sum_{k=n+1}^{\infty} h[k]\,e^{-j\omega k}\,e^{j\omega n}\right| \\
&\le \sum_{k=n+1}^{\infty} |h[k]|\,\left|e^{-j\omega k}\,e^{j\omega n}\right| \\
&\le \sum_{k=n+1}^{\infty} |h[k]| \le \sum_{k=0}^{\infty} |h[k]|
\end{aligned} \]

\sayfa{5}
Sistem kararlı ise transient'i de kararlıdır.

\section{3. bölüm}

\[ X(e^{j\omega}) = \sum_{n=-\infty}^{\infty} x[n]\,e^{-j\omega n} \]
\[ X(z) = \sum_{n=-\infty}^{\infty} x[n]\,z^{-n} \]

\[ z = re^{j\omega}\ \Big| \]
\[ \begin{aligned}
&= \sum_{n=-\infty}^{\infty} x[n]\,(re^{j\omega})^{-n} \\
&= \sum_{n=-\infty}^{\infty} \left[x[n]\,r^{-n}\right] e^{-j\omega n}
\end{aligned} \]

z dönüşümü $x[n]r^{-n}$'in Fourier Transform \\
kararlılık analizi için kullanılır \\
$z = re^{j\omega}$

\elcizimi[0.35]{sekiller/Dsp-ders-4-s5-1.png}[5]

\elcizimi[0.4]{sekiller/Dsp-ders-4-s5-2.png}[5]

$r=1$ için z transformu ile \okunamadi{Fourier} dönüşümü aynıdır

\elcizimi[0.6]{sekiller/Dsp-ders-4-s5-3.png}[5]

Sonuç, birim çemberin üzerindeyse sistem kararlıdır

\[ \begin{aligned} x[n] &\donusum X(z) \\ h[n] &\donusum H(z) \\ y[n] &\donusum Y(z) \end{aligned} \]

\[ X(z) = \frac{P(z)}{Q(z)} \quad \text{(rasyonel)} \]

$P(z)$'nin kökleri $X(z)$ sıfırı \\
$Q(z)$'nin kökleri kutup

\sayfa{6}
\begin{ornek}[1]
$x[n]=a^n u[n]$, \quad $X(z)=\,?$
\end{ornek}

\[ \begin{aligned}
X(z) &= \sum_{n=-\infty}^{\infty} a^n u[n]\cdot z^{-n} \\
&= \sum_{n=-\infty}^{\infty} \left(\frac{a}{z}\right)^{n} u[n] \\
&= \frac{1}{1-\frac{a}{z}} \qquad \left|\frac{a}{z}\right|<1 \qquad |z|>|a|
\end{aligned} \]
{\small $\displaystyle\sum_{n=0}^{\infty} a^n$}

\[ X(z) = \frac{1}{1-az^{-1}} \]

\[ \begin{aligned}
X(z) &= \frac{1}{1-az^{-1}} \ \to\ \text{negatif form} \\
&= \frac{z}{z-a}
\end{aligned} \]

\begin{dikkat}
\okunamadi{mutlaka} bu forma getir.
\end{dikkat}

sıfırda bir sıfırı var \okunamadi{baştan göremedik}

\elcizimi[0.65]{sekiller/Dsp-ders-4-s6-1.png}[6]

\begin{ornek}[2]
$x[n] = -a^n u[-n-1]$
\end{ornek}

\[ \begin{aligned}
X(z) &= \sum_{n=-\infty}^{\infty} -a^n u[-n-1]\,z^{-n} \\
&= \sum_{n=-\infty}^{-1} -\left(\frac{a}{z}\right)^{n} = \sum_{n=1}^{\infty} -\left(\frac{a}{z}\right)^{-n} = \sum_{n=0}^{\infty} -\left(\frac{a}{z}\right)^{-n} - [-1]
\end{aligned} \]

\formulkutu[önemli]{\sum_{n=0}^{\infty} a^n = \frac{1}{1-a} \qquad |a|<1}

\[ \begin{aligned}
&= 1 - \sum_{n=0}^{\infty} \left(\frac{z}{a}\right)^{n} \\
&= 1 - \frac{1}{1-\frac{z}{a}} \qquad \begin{matrix} |a|<1 \\ \left|\frac{z}{a}\right|<1 \end{matrix} \\
&= 1 - \frac{a}{a-z} = \frac{-z}{a-z} \\
&= \frac{z}{z-a} \qquad \begin{matrix} \left|\frac{z}{a}\right|<1 \\ |z|<a \end{matrix}
\end{aligned} \]

z dönüşümü ROC'la birlikte anlamlıdır

\elcizimi[0.55]{sekiller/Dsp-ders-4-s6-2.png}[6]

\sayfa{7}
\begin{ornek}[3]
$x[n] = \left(\tfrac12\right)^{n} u[n] + \left(-\tfrac13\right)^{n} u[n]$
\end{ornek}

\[ X(z) = \frac{z}{z-\frac12} + \frac{z}{z+\frac13} = \frac{2z\left(z-\frac{1}{12}\right)}{\left(z-\frac12\right)\left(z+\frac13\right)} \qquad |z|>\tfrac12 \]

\elcizimi[0.9]{sekiller/Dsp-ders-4-s7-1.png}[7]

\begin{ornek}[4]
$x[n] = \left(-\tfrac13\right)^{n} u[n] - \left(\tfrac12\right)^{n} u[-n-1]$
\end{ornek}

\[ X(z) = \underset{|z|>\frac13}{\frac{z}{z+\frac13}} + \underset{|z|<\frac12}{\frac{z}{z-\frac12}} = \underset{\frac13<|z|<\frac12}{\frac{2z\left(z-\frac{1}{12}\right)}{\left(z-\frac12\right)\left(z+\frac13\right)}} \]

\elcizimi[0.5]{sekiller/Dsp-ders-4-s7-2.png}[7]

\sayfa{8}
\begin{ornek}[5]
$x[n] = \left(-\tfrac12\right)^{n} u[n] - \left(\tfrac13\right)^{n} u[-n-1]$
\end{ornek}

\[ X(z) = \underset{|z|>\frac12}{\frac{z}{z+\frac12}} + \underset{|z|<\frac13}{\frac{z}{z-\frac13}} \]

ROC yok \\
hiç bir zaman çakışmaz

\begin{ornek}[6]
\[ x[n] = \begin{cases} a^{n}, & 0\leq n\leq N-1 \\ 0, & \text{o.w.} \end{cases} \]
\end{ornek}

\[ \begin{aligned}
X(z) &= \sum_{0}^{N-1} a^n z^{-n} = 1 + az^{-1} + a^2 z^{-2} + \cdots + a^{N-1} z^{-(N-1)} \\
&= 1 + \frac{a}{z} + \frac{a^2}{z^2} + \cdots + \frac{a^{N-1}}{z^{N-1}}
\end{aligned} \]

$z=0$'da toplam sonsuz

z hariç her yerde z dönüşümü vardır

\[ \frac{1-a^{N}}{1-az^{-1}}\footnotemark \]
\footnotetext{Claude notu: muhtemelen pay $1-a^{N}z^{-N}$ olmalı.}

\begin{ornek}[7]
\[ x[n] = \begin{cases} a^{n}, & -N\leq n\leq 0 \\ 0, & \text{o.w} \end{cases} \]
\end{ornek}

\[ X(z) = 1 + a^{-1}z + a^{-2}z^{2} + \cdots + a^{-N}z^{N} \]

$z=\infty$ için $\infty$ olur

$z\neq\infty$ için z dönüşümü vardır

\elcizimi[0.45]{sekiller/Dsp-ders-4-s8-1.png}[8]

\[ H(z) = K\,\frac{z}{\left(z-\frac12\right)(z+2)} \]
\okunamadi{Köksal'ın} K'sı :)

\[ |z|>2 \quad \checkmark \to \text{causal} \]
\[ \tfrac12<|z|<2 \quad \checkmark\ \text{stable.} \]
\[ |z|<\tfrac12 \quad \text{ne nedensel ne kararlıdır} \]

kararlı ise ROC birim çemberi içermeli

$\star$ Sorunun \okunamadi{içinde} nedensel derse sağ yanlıdır

\end{document}
